\documentclass[11pt]{article}

%%% Packages
\usepackage{amsmath,amsthm,amssymb,amsfonts}
\usepackage{mathtools}
\usepackage{enumerate}
\usepackage{hyperref}
\usepackage{url}
\usepackage{booktabs}
\usepackage{array}
\usepackage{xcolor}
\usepackage[margin=1in]{geometry}

%%% Theorem environments
\newtheorem{theorem}{Theorem}[section]
\newtheorem{lemma}[theorem]{Lemma}
\newtheorem{proposition}[theorem]{Proposition}
\newtheorem{corollary}[theorem]{Corollary}
\newtheorem{conjecture}[theorem]{Conjecture}
\theoremstyle{definition}
\newtheorem{definition}[theorem]{Definition}
\newtheorem{example}[theorem]{Example}
\newtheorem{remark}[theorem]{Remark}
\newtheorem{observation}[theorem]{Observation}

%%% Notation shortcuts
\newcommand{\Z}{\mathbb{Z}}
\newcommand{\Q}{\mathbb{Q}}
\newcommand{\R}{\mathbb{R}}
\newcommand{\C}{\mathbb{C}}
\newcommand{\N}{\mathbb{N}}
\newcommand{\F}{\mathcal{F}}
\newcommand{\fpart}[1]{\left\{#1\right\}}
\DeclareMathOperator{\rank}{rank}
\DeclareMathOperator{\sgn}{sgn}
\DeclareMathOperator{\Res}{Res}
\DeclareMathOperator{\Cov}{Cov}
\DeclareMathOperator{\Var}{Var}

\title{Per-Step Farey Discrepancy, Zeta Zero Oscillations,\\
and a Chebyshev Bias for Rational Number Regularity}

\author{
  Saar Shai\thanks{Independent researcher. Email: \texttt{saar.shai@gmail.com}.}
  \and
  Claude Opus 4.6 (Anthropic)\thanks{AI research assistant; contributed to mathematical analysis, computation, and manuscript preparation.}
}

\date{March 2026}

\begin{document}

\maketitle

\begin{abstract}
Per-step Farey discrepancy $\Delta W(N) = W(N-1) - W(N)$ has never been
studied.  It reveals structure invisible to the classical cumulative
discrepancy~$W(N)$, and led to the following discoveries:
\begin{enumerate}[(i)]
  \item An exact four-term algebraic decomposition of $\Delta W(p)$ for
    every prime~$p$ (Theorem~\ref{thm:four-term}), reducing the sign
    of~$\Delta W$ to an explicit number-theoretic sum $T(N)$.
  \item A regression identity $B'/C' = \alpha + \rho$ (Theorem~\ref{thm:regression}),
    proved via permutation invariance of coprime residues.
  \item Computational verification that $\Delta W(p) < 0$ for all
    4{,}617 primes $p \le 10^5$ with $M(p) \le -3$, and discovery of the
    first counterexample at $p = 243{,}799$.
  \item A Perron integral formula for $T(N)$ (Theorem~\ref{thm:perron},
    conditional on GRH), expressing $T(N)$ as $-2\log N + C_0
    + \sum_\rho c_\rho N^\rho + O(1)$ with computable residues~$c_\rho$.
  \item A phase-lock theorem and Chebyshev bias
    (Theorem~\ref{thm:phase-lock}, conditional on GRH+LI):
    the sign of $\Delta W(p)$ is controlled by
    $\gamma_1 \log p \bmod 2\pi$, where $\gamma_1 = 14.134\ldots$ is the
    first zeta zero.  The predicted phase matches observation to
    $0.07$~radians.
  \item A spectral decomposition indexed by Dirichlet characters
    (Section~\ref{sec:l-function}): each twisted summatory function~$T_\chi$
    oscillates at frequencies determined by zeros of~$L(s,\chi)$.
  \item A null test (Section~\ref{sec:ordering}): the Stern--Brocot tree
    enumerates the same rationals but produces no spectral signature,
    confirming that the zeta-zero oscillation is a property of the
    \emph{arithmetic ordering by denominator}, not of the rationals
    themselves.
  \item Unconditional results (Theorem~\ref{thm:unconditional}):
    $T(N) = \Omega_\pm(\sqrt{N})$ via Ingham's method, giving
    infinitely many primes with $\Delta W > 0$ and infinitely many
    with $\Delta W < 0$.
\end{enumerate}
\end{abstract}

\medskip
\noindent\textbf{Keywords:} Farey sequence, per-step discrepancy,
Mertens function, Chebyshev bias, zeta zeros, Dirichlet $L$-functions,
Stern--Brocot tree.

\noindent\textbf{MSC 2020:} 11B57, 11N37, 11M06, 11M26, 11A25.



%% ================================================================
%% SECTION 1: INTRODUCTION
%% ================================================================
\section{Introduction}\label{sec:intro}

\subsection{Farey sequences and the Franel--Landau theorem}

For a positive integer~$N$, the \emph{Farey sequence} $\F_N$ is the
ascending list of fractions $a/b$ with $0 \le a \le b \le N$ and
$\gcd(a,b) = 1$.  Its cardinality is
\begin{equation}\label{eq:farey-size}
  |\F_N| = 1 + \sum_{k=1}^{N} \varphi(k)
  = \frac{3N^2}{\pi^2} + O(N \log N).
\end{equation}
Writing $n = |\F_N|$ and listing the elements as
$f_0 < f_1 < \cdots < f_{n-1}$, the \emph{$L^2$ discrepancy} is
\begin{equation}\label{eq:wobble}
  W(N) = \frac{1}{n^2} \sum_{j=0}^{n-1}
    \Bigl(f_j - \frac{j}{n}\Bigr)^{\!2}.
\end{equation}
Franel~\cite{Franel1924} and Landau~\cite{Landau1924} proved that the
Riemann Hypothesis is equivalent to $W(N) = O(N^{-1+\varepsilon})$ for
every~$\varepsilon > 0$.

\subsection{The per-step question: a new perspective}

Despite decades of work on the Franel--Landau framework, the
\emph{per-step} behavior---how $W$ changes when a single integer
enters---appears not to have been studied.  Define
\begin{equation}\label{eq:delta-w}
  \Delta W(N) = W(N{-}1) - W(N),
\end{equation}
so $\Delta W(N) > 0$ means adding $N$ \emph{improves} uniformity.
The telescoping identity $W(N) = W(1) - \sum_{k=2}^{N} \Delta W(k)$
decomposes the cumulative discrepancy into individual contributions.

The per-step viewpoint transforms a static asymptotic question
(how fast does $W(N) \to 0$?) into a dynamical one (which integers
improve the distribution, and which degrade it?).  This dynamical
perspective reveals oscillatory structure locked to the nontrivial
zeros of the Riemann zeta function---structure that is invisible in the
cumulative discrepancy.

%% FIGURE: DeltaW(N) time series for N = 2, ..., 1000 showing oscillation
%% between positive (composites healing) and negative (primes damaging) values.
%% Mark the prime contributions in a different color.

\subsection{Summary of discoveries}\label{sec:intro-results}

The per-step framework yields a chain of results connecting elementary
algebra to the deepest structure in analytic number theory.

\begin{enumerate}[(1)]
  \item \textbf{Four-term decomposition} (Theorem~\ref{thm:four-term}).
    An exact algebraic identity
    $n'^2 \Delta W(p) = A - B - C - D$,
    decomposing each prime step into four arithmetically interpretable
    terms.  Moreover, $|D/A - 1| = O(1/p^2)$, so the sign of
    $\Delta W(p)$ is controlled by $B + C$.

  \item \textbf{Regression identity} (Theorem~\ref{thm:regression}).
    $B'/C' = \alpha + \rho$,
    where $\alpha$ is a regression slope connected to the Mertens
    function and $\rho \approx -3.9$ is a stable residual.  This
    is an exact algebraic identity for every prime $p \ge 5$.

  \item \textbf{Computational Sign Theorem} (Section~\ref{sec:sign}).
    $\Delta W(p) < 0$ for all 4{,}617 primes $p \le 10^5$ with
    $M(p) \le -3$, with zero exceptions.  The first counterexample
    occurs at $p = 243{,}799$, where $T(N)$ crosses zero for the
    first time on this subsequence.

  \item \textbf{Perron integral for $T(N)$}
    (Theorem~\ref{thm:perron}, conditional on GRH).
    \[
      T(N) = -2\log N + C_0 + \sum_\rho c_\rho N^\rho + O(1),
    \]
    with $c_\rho = \zeta(\rho+1)/(\rho\,\zeta'(\rho))$ and
    $C_0 = 2 - 2\gamma + 2\log 2\pi = 4.521\ldots$

  \item \textbf{Phase-lock theorem and Chebyshev bias}
    (Theorem~\ref{thm:phase-lock}, conditional on GRH+LI).
    The logarithmic density of $\{p : \Delta W(p) > 0\}$
    among $M(p) = -3$ primes is~$1/2$.  The sign of $\Delta W(p)$
    is controlled by the phase $\gamma_1 \log p \bmod 2\pi$.
    The predicted phase peak is $5.208$; observed: $5.28$
    (error $0.07$ radians, $1.1\%$).

  \item \textbf{L-function spectral decomposition}
    (Section~\ref{sec:l-function}).
    Per-step Farey discrepancy admits a spectral decomposition indexed
    by Dirichlet characters, where each component oscillates at
    frequencies determined by the corresponding $L$-function's
    nontrivial zeros.

  \item \textbf{Ordering principle} (Section~\ref{sec:ordering}).
    Different orderings of the same rationals reveal different
    spectral data---or none at all.  The Stern--Brocot tree produces
    monotone exponential growth in $\Delta W$ with ratio $\to 2$
    and zero spectral content.  The Farey ordering by denominator
    is special because it engages multiplicative number theory,
    making zeta zeros visible.

  \item \textbf{Unconditional oscillation}
    (Theorem~\ref{thm:unconditional}).
    $T(N) = \Omega_\pm(\sqrt{N})$, proved by adapting Ingham's 1942
    method.  Both $\Delta W > 0$ and $\Delta W < 0$ occur infinitely
    often among $M(p) = -3$ primes.
\end{enumerate}

\subsection{What this framework does not do}\label{sec:intro-honest}

The individual techniques used in this paper are standard: the Perron
integral (Titchmarsh, Ch.~14), the Rubinstein--Sarnak framework~\cite{RubinsteinSarnak1994},
Farey sequence asymptotics (Franel 1924, Landau 1924), and OLS
decomposition.  What is new is the \emph{per-step perspective} and the
specific results it yields.  In particular:

\begin{itemize}
  \item This is \textbf{not} a new route to the Riemann Hypothesis.
    The per-step sign controls the \emph{direction} of change, not
    the \emph{magnitude}.
  \item The Sign Theorem is \emph{not} universal: it fails at
    $p = 243{,}799$, and the fraction of counterexamples grows
    toward~$1/2$.
  \item The phase-lock and density-$1/2$ results are
    \emph{conditional} on GRH and the Linear Independence hypothesis.
\end{itemize}

\subsection{Paper outline}

Section~\ref{sec:defs} establishes definitions and the four-term
decomposition.
Section~\ref{sec:regression} proves the regression identity.
Section~\ref{sec:sign} presents the computational Sign Theorem and
the counterexample mechanism.
Section~\ref{sec:perron} derives the Perron integral for~$T(N)$.
Section~\ref{sec:phaselock} proves the phase-lock theorem and
Chebyshev bias.
Section~\ref{sec:l-function} develops the $L$-function spectral
decomposition.
Section~\ref{sec:ordering} presents the Stern--Brocot null test.
Section~\ref{sec:unconditional} gives unconditional results via
Ingham's method.
Section~\ref{sec:multiscale} analyzes the multi-scale structure
of~$T(N)$.
Section~\ref{sec:open} lists open problems.



%% ================================================================
%% SECTION 2: DEFINITIONS AND FOUR-TERM DECOMPOSITION
%% ================================================================
\section{Definitions and the Four-Term Decomposition}\label{sec:defs}

\subsection{Setup}

Throughout, let $p \ge 5$ be prime, $N = p - 1$, $n = |\F_N|$,
$n' = |\F_p| = n + p - 1$.

\begin{definition}[Rank discrepancy]
For $f_j \in \F_N$, the rank discrepancy is
$D(f_j) = j - n \cdot f_j$.
\end{definition}

\begin{definition}[Per-step shift]
For $f = a/b \in \F_{N}$ with $b \ge 2$,
\begin{equation}\label{eq:shift}
  \delta(f) = \frac{a - (pa \bmod b)}{b}.
\end{equation}
Since $\gcd(p,b) = 1$, the map $a \mapsto pa \bmod b$ is a
permutation of the residues coprime to~$b$.
\end{definition}

\begin{definition}[Regression quantities]
Let $F$ denote the interior fractions of $\F_N$ (denominators $\ge 2$).
Define:
\begin{itemize}
  \item $\bar{f} = \frac{1}{|F|}\sum_{f \in F} f = \frac{1}{2}$
    (Farey symmetry);
  \item $\bar{D} = \frac{1}{|F|}\sum_{f \in F} D(f) = \frac{1}{2}$;
  \item $\alpha = \Cov(D, f)/\Var(f)$, the OLS slope;
  \item $D_{\mathrm{err}}(f) = D(f) - \bar{D} - \alpha(f - \tfrac{1}{2})$,
    the OLS residual;
  \item $B' = 2\sum_{f \in F} D(f)\,\delta(f)$;
  \item $C' = \sum_{f \in F} \delta(f)^2$;
  \item $\rho = 2\sum_{f \in F} D_{\mathrm{err}}(f)\,\delta(f) \,/\, C'$.
\end{itemize}
\end{definition}

\begin{definition}[The sum $T(N)$]
\begin{equation}\label{eq:T-def}
  T(N) = \sum_{m=2}^{N} \frac{M(\lfloor N/m \rfloor)}{m},
\end{equation}
where $M(x) = \sum_{k \le x} \mu(k)$ is the Mertens function.
\end{definition}

\subsection{The displacement--shift identity}

\begin{proposition}[Displacement--shift identity]\label{prop:disp-shift}
For prime~$p$ and $f = a/b \in \F_{p-1}$ with $f \ne 1$:
\[
  D_{\F_p}(f) = D_{\F_{p-1}}(f) + \delta(f).
\]
\end{proposition}

\begin{proof}
The rank of $f$ in $\F_p$ equals its rank in $\F_{p-1}$ plus the
number of new fractions $k/p < f$, which is $\lfloor pa/b \rfloor$.
Then:
\begin{align*}
  D_{\F_p}(f) &= D_{\F_{p-1}}(f) + \lfloor pa/b \rfloor - (p-1)f \\
  &= D_{\F_{p-1}}(f) + \bigl(pa/b - (pa \bmod b)/b\bigr) - (p-1)(a/b) \\
  &= D_{\F_{p-1}}(f) + \delta(f). \qedhere
\end{align*}
\end{proof}

\subsection{The four-term decomposition}

\begin{theorem}[Four-term decomposition]\label{thm:four-term}
For every prime $p \ge 5$:
\begin{equation}\label{eq:four-term}
  n'^2 \Delta W(p) = A - B - C - D,
\end{equation}
where:
\begin{align}
  A &= \Bigl(\frac{n'^2}{n^2} - 1\Bigr)
    \sum_{f \in \F_N} D_{\F_N}(f)^2
    && \text{(dilution gain),} \\
  B &= 2\sum_{f \in F} D_{\F_N}(f)\,\delta(f)
    && \text{(cross term),} \\
  C &= \sum_{f \in F} \delta(f)^2
    && \text{(shift-squared cost),} \\
  D &= \sum_{k=1}^{p-1} D_{\F_p}(k/p)^2
    && \text{(new-fraction cost).}
\end{align}
Moreover, $|D/A - 1| = O(1/p^2)$ as $p \to \infty$.
\end{theorem}

\begin{proof}
Expanding $n'^2 W(p) = \sum_{f \in \F_p} D_{\F_p}(f)^2$ via the
displacement--shift identity:
\begin{align*}
  n'^2 W(p) &= \sum_{f \in \F_N} (D_{\F_N}(f) + \delta(f))^2
    + \sum_{k=1}^{p-1} D_{\F_p}(k/p)^2 \\
  &= n^2 W(N) + B + C + D.
\end{align*}
Then $n'^2 \Delta W(p) = n'^2 W(N) - n'^2 W(p)
= (n'^2 - n^2)W(N) - B - C - D = A - B - C - D$.

For $D/A \to 1$: writing the new-fraction ranks as
$D_{\F_p}(k/p) = k - n'(k/p)$, one has $D = \Theta(p^3)$ and
$A = \Theta(p^3)$, with $|D/A - 1| = O(1/p^2)$ from matching the
leading terms.
\end{proof}

Since $D \approx A$, the sign of $\Delta W(p)$ is controlled by
$B + C$: positive $B + C$ means the cross and shift terms absorb
the dilution-cost difference, giving $\Delta W < 0$ (regularity
worsens).



%% ================================================================
%% SECTION 3: THE REGRESSION IDENTITY
%% ================================================================
\section{The Regression Identity $B'/C' = \alpha + \rho$}
\label{sec:regression}

\begin{lemma}[Shift residual vanishing]\label{lem:shift-vanish}
For each denominator $b \in \{2, \ldots, N\}$:
\[
  \sum_{\substack{a=1 \\ \gcd(a,b)=1}}^{b-1} \delta(a/b) = 0.
\]
\end{lemma}

\begin{proof}
Since $p$ is prime and $b < p$, the map $a \mapsto pa \bmod b$ is a
permutation of the coprime residues.  Hence
$\sum_a \delta(a/b) = \frac{1}{b}[\sum_a a - \sum_a a] = 0$.
\end{proof}

\begin{corollary}\label{cor:sum-delta-zero}
$\sum_{f \in F} \delta(f) = 0$.
\end{corollary}

\begin{lemma}[Permutation square-sum identity]\label{lem:perm-sq}
$\sum_{f \in F} f \cdot \delta(f) = C'/2$.
\end{lemma}

\begin{proof}
Write $\delta(f) = f - \{pf\}$.  Then
$f\,\delta(f) - \delta(f)^2/2 = (f^2 - \{pf\}^2)/2$.
For each fixed~$b$, the permutation $a \mapsto pa \bmod b$ gives
$\sum_a (a/b)^2 = \sum_a ((pa \bmod b)/b)^2$.  Summing over all~$b$
yields $\sum f\,\delta = C'/2$.
\end{proof}

\begin{theorem}[Regression identity]\label{thm:regression}
For every prime $p \ge 5$:
\begin{equation}\label{eq:regression}
  \frac{B'}{C'} = \alpha + \rho,
\end{equation}
where $\alpha = \Cov(D,f)/\Var(f)$ and
$\rho = 2\sum D_{\mathrm{err}}(f)\,\delta(f)/C'$.
\end{theorem}

\begin{proof}
Decompose $D(f) = \bar{D} + \alpha(f - 1/2) + D_{\mathrm{err}}(f)$.  Then:
\[
  B' = 2\bar{D}\underbrace{\sum\delta}_{=\,0}
  + 2\alpha\underbrace{\sum(f - \tfrac{1}{2})\delta}_{=\,C'/2}
  + 2\sum D_{\mathrm{err}}\,\delta
  = \alpha C' + \rho C'.  \qedhere
\]
\end{proof}

\begin{remark}
This identity is purely algebraic: no approximations, no asymptotics.
It holds exactly for every prime $p \ge 5$.  The connection to~$T(N)$
comes through the classical relation $\alpha \approx 1 - T(N)$ (see
Theorem~\ref{thm:alpha-formula}), which links the OLS slope to the
Mertens function at multiple scales.
\end{remark}

\begin{theorem}[Alpha formula]\label{thm:alpha-formula}
For prime~$p$ with $M(p-1) = -2$:
\[
  \alpha = 1 - T(N) + O(1/N),
\]
where $T(N) = \sum_{m=2}^{N} M(\lfloor N/m \rfloor)/m$ and $N = p-1$.
\end{theorem}

\begin{proof}[Proof sketch]
The regression slope $\alpha = \Cov(D,f)/\Var(f)$ involves the quantity
$\sum f \cdot D(f)$, which by classical Farey asymptotics (Franel--Landau)
equals $-nR(N)/2 + O(1)$, where
$R(N) = \frac{1}{6} + \frac{1}{6}\sum_{m=1}^{N} M(\lfloor N/m\rfloor)/m$.
Since $\Var(f) = 1/12 + O(1/N)$, we get $\alpha = -6R(N) + O(1/N)$.
Expanding: $6R(N) = 1 + M(N) + T(N)$, and for $M(N) = -2$:
$\alpha = 1 - T(N) + O(1/N)$.
\end{proof}

The chain of identities is now complete:
\[
  \Delta W(p) < 0
  \;\Longleftrightarrow\; B + C > 0
  \;\Longleftrightarrow\; \alpha + \rho > -1
  \;\Longleftrightarrow\; T(N) < 2 + \rho.
\]
Since $\rho \approx -3.9$ for large~$p$, the threshold is approximately
$T(N) < -1.9$.



%% ================================================================
%% SECTION 4: COMPUTATIONAL SIGN THEOREM AND COUNTEREXAMPLE
%% ================================================================
\section{Computational Sign Theorem and the First Counterexample}
\label{sec:sign}

\subsection{Sign Theorem to $10^5$}

\begin{theorem}[Computational]\label{thm:sign}
$\Delta W(p) < 0$ for all $4{,}617$ primes $11 \le p \le 10^5$ with
$M(p) \le -3$.
\end{theorem}

This was verified by exact rational arithmetic over all qualifying
primes.  For each prime, the four-term decomposition was evaluated
by enumerating $\F_{p-1}$ via the mediant algorithm and computing
$A$, $B$, $C$, $D$ exactly, with internal consistency checks
($\sum \delta = 0$ per denominator, $\sum f\delta = C'/2$).

The mechanism: for $p \le 10^5$, the sum $T(N)$ is uniformly negative
(reaching $-50$ at $N \approx 10^5$), so $\alpha \approx 1 - T(N) \approx 51$,
which far exceeds $|\rho| \approx 3.9$.  The inequality
$\alpha + \rho > -1$ holds with enormous margin.

\subsection{The first counterexample: $p = 243{,}799$}

\begin{proposition}[Computational]\label{prop:counterexample}
The first prime $p$ with $M(p) \le -3$ and $\Delta W(p) > 0$ is
$p = 243{,}799$.
\end{proposition}

At this prime, $N = 243{,}798$, $M(N) = -2$, and:

\begin{center}
\begin{tabular}{ll}
\toprule
Quantity & Value \\
\midrule
$T(N)$ & $+0.165$ \\
$\alpha$ & $0.835$ \\
$\rho$ & $-3.887$ \\
$B'/C'$ & $-3.052$ \\
$D/A$ & $0.981$ \\
$n'^2 \Delta W$ & $+6.65 \times 10^9$ \\
\bottomrule
\end{tabular}
\end{center}

\textbf{Failure mechanism.}  $T(N)$ crosses zero for the first time on
the $M(N) = -2$ subsequence, causing $\alpha$ to drop from its typical
value of $\sim 12$ (at $p \sim 10^5$) to~$0.835$.  Since
$\rho \approx -3.9$ is stable, $\alpha + \rho = -3.05 < -1$, and
$\Delta W > 0$.

\textbf{This counterexample is not an anomaly}---it is the onset of
a systematic pattern.  Among the 922 primes $p \le 10^7$ with
$M(p) = -3$, 247 (26.8\%) have $T(N) > 0$ and hence $\Delta W > 0$.
The fraction increases with~$p$: from $0\%$ below $10^5$, to $10\%$
in $[10^5, 10^6)$, to $46\%$ in $[10^6, 10^7)$.

%% FIGURE: Running fraction of T(N) > 0 primes vs p, showing the
%% monotone increase from 0 toward ~1/2.  Include the prediction
%% from the Rubinstein-Sarnak framework (dashed curve) for comparison.

\subsection{Rho stability}

The residual $\rho$ stabilizes rapidly:

\begin{center}
\begin{tabular}{ll}
\toprule
Range of $p$ & $\rho$ range \\
\midrule
$[13, 179]$ & $[-3.2, -1.3]$ \\
$[1000, 5000]$ & $[-3.85, -3.64]$ \\
$[5000, 100{,}000]$ & $[-3.89, -3.83]$ \\
\bottomrule
\end{tabular}
\end{center}

For $p > 5000$, $\rho$ is effectively constant at $\approx -3.88$.
The sign of $\Delta W$ is therefore determined by $\alpha$, hence
by~$T(N)$.



%% ================================================================
%% SECTION 5: THE PERRON INTEGRAL FOR T(N)
%% ================================================================
\section{The Perron Integral for $T(N)$}\label{sec:perron}

\begin{theorem}[Conditional on GRH]\label{thm:perron}
Assume the Generalized Riemann Hypothesis.  Then for $M(N) = -2$:
\begin{equation}\label{eq:perron}
  T(N) = -2\log N + C_0 + \sum_{\rho} c_\rho\, N^\rho + O(1),
\end{equation}
where $C_0 = 2(1 - \gamma + \log 2\pi) = 4.521\ldots$,
$c_\rho = \zeta(\rho+1)/(\rho\,\zeta'(\rho))$,
and the sum is over nontrivial zeros~$\rho$ of $\zeta(s)$.
\end{theorem}

\begin{proof}
Define $F(N) = \sum_{m=1}^{N} M(\lfloor N/m \rfloor)/m$, so
$T(N) = F(N) - M(N)$.  The Perron integral representation is
\begin{equation}\label{eq:perron-integral}
  F(N) = \frac{1}{2\pi i} \int_{(c)} N^s \cdot
  \frac{\zeta(s+1)}{s\,\zeta(s)}\,ds, \qquad c > 1,
\end{equation}
using the Dirichlet series identities: $\mu(n)$ has generating
function $1/\zeta(s)$, the harmonic weight $1/n$ has generating
function $\zeta(s+1)$, and the Perron formula
$\frac{1}{2\pi i}\int_{(c)} x^s/s\,ds = \mathbf{1}_{x > 1}$.

The integrand $G(s) = N^s\,\zeta(s+1)/(s\,\zeta(s))$ has:
\begin{enumerate}[(i)]
  \item A double pole at $s = 0$ (from $1/s$ and the pole of
    $\zeta(s+1)$);
  \item Simple poles at $s = \rho$ for each nontrivial zero of
    $\zeta(s)$;
  \item Trivial-zero poles at $s = -2k$ ($k = 1, 2, \ldots$),
    contributing $O(1)$.
\end{enumerate}

\textbf{Residue at $s = 0$.}  Using the Laurent expansions
$N^s = 1 + s\log N + \cdots$,
$\zeta(s+1) = 1/s + \gamma + \cdots$,
$1/\zeta(s) = -2 + 2\log(2\pi)\cdot s + \cdots$
(since $\zeta(0) = -1/2$):
\[
  \Res_{s=0} G(s) = -2(\log N + \gamma) + 2\log 2\pi
  = -2\log N + 2(\log 2\pi - \gamma).
\]

\textbf{Residue at $s = \rho$.}
$\Res_{s=\rho} G(s) = N^\rho \cdot \zeta(\rho+1)/(\rho\,\zeta'(\rho))
= c_\rho\, N^\rho$.

Shifting the contour to $\operatorname{Re}(s) = -1 - \varepsilon$
(justified under GRH):
\[
  F(N) = -2\log N + 2(\log 2\pi - \gamma)
  + \sum_\rho c_\rho\, N^\rho + O(1).
\]
Since $T(N) = F(N) - M(N)$ and $M(N) = -2$, we obtain~\eqref{eq:perron}.
\end{proof}

\subsection{Computed coefficients}

Writing $\rho_k = 1/2 + i\gamma_k$ (assuming RH):

\begin{center}
\begin{tabular}{ccccc}
\toprule
$k$ & $\gamma_k$ & $|c_k|$ & $\arg(c_k)$ & Period $2\pi/\gamma_k$ \\
\midrule
1 & 14.1347 & 0.04853 & $-1.602$ & 0.445 \\
2 & 21.0220 & 0.02778 & $-1.487$ & 0.299 \\
3 & 25.0109 & 0.02170 & $-1.668$ & 0.251 \\
4 & 30.4249 & 0.01704 & $-1.337$ & 0.207 \\
5 & 32.9351 & 0.01526 & $-1.799$ & 0.191 \\
\bottomrule
\end{tabular}
\end{center}

Computed with \texttt{mpmath} at 30-digit precision.  The dominant
oscillatory term is:
\begin{equation}\label{eq:dominant-osc}
  T(N) \approx -2\log N + 4.52
  + 0.0971\sqrt{N}\cos(14.1347\log N - 1.602) + \cdots
\end{equation}
The oscillatory amplitude $O(\sqrt{N})$ eventually dominates the
drift $-2\log N$, explaining why counterexamples emerge.

%% FIGURE: Overlay of the Perron prediction (first 5 zeros) vs
%% observed T(N) values for 922 M(p)=-3 primes.  The agreement
%% should be visible in both the overall trend and the oscillation
%% phases.



%% ================================================================
%% SECTION 6: PHASE-LOCK THEOREM AND CHEBYSHEV BIAS
%% ================================================================
\section{Phase-Lock Theorem and Chebyshev Bias}\label{sec:phaselock}

\begin{theorem}[Phase-lock and Chebyshev bias; conditional on GRH+LI]
\label{thm:phase-lock}
Assume GRH and the Linear Independence hypothesis
($\gamma_1, \gamma_2, \ldots$ are linearly independent over~$\Q$).
Then:

\begin{enumerate}[\upshape(a)]
  \item The logarithmic density
  \[
    \delta = \lim_{X \to \infty} \frac{1}{\log X}
    \int_2^X \mathbf{1}_{T(t) > 0}\,\frac{dt}{t}
  \]
  exists and equals~$1/2$.

  \item Among $N$ with $T(N) > 0$, the circular mean of
  $\theta(N) = \gamma_1 \log N \bmod 2\pi$ concentrates at the
  phase $\gamma_1 \log 2 - \arg(d_1) \pmod{2\pi}$, where
  $d_1 = 1/(\rho_1\,\zeta'(\rho_1))$.

  \item The finite-$N$ bias is $O(\log N / \sqrt{N})$.
\end{enumerate}
\end{theorem}

\begin{proof}
By Theorem~\ref{thm:perron}, $T(N)/\sqrt{N}$ converges in distribution
(under logarithmic measure) to
\[
  Y = \sum_{\gamma_k > 0} 2|c_k|\cos(\theta_k + \arg c_k),
\]
where $\theta_k$ are independent uniform on $[0, 2\pi)$ (this is the
Rubinstein--Sarnak framework~\cite{RubinsteinSarnak1994}, made rigorous
under GRH+LI).

\textbf{(a)} Each cosine term is symmetric about zero, so $Y$ has a
symmetric distribution.  The condition $T(N) > 0$ becomes
$Y > (2\log N - C_0)/\sqrt{N} \to 0$, so the density equals
$\Pr(Y > 0) = 1/2$.

\textbf{(b)} The $M(N/2)/2$ term dominates $T(N)$ (empirical
correlation $r = 0.95$).  The Mertens function oscillates as
$M(x) \sim 2\operatorname{Re}[d_1 x^{\rho_1}]$, where
$|d_1| = 0.0891$, $\arg(d_1) = -1.693$.  Setting $x = N/2$:
$M(N/2) > 0$ when
$\gamma_1 \log N \approx \gamma_1 \log 2 - \arg(d_1) = 5.208$.

\textbf{(c)} The effective threshold
$\delta(N) = (2\log N - C_0)/\sqrt{N}$ gives
$\Pr(T > 0) = 1/2 - O(\delta(N)/\sigma)$ where
$\sigma^2 = \sum_k 2|c_k|^2 = 0.01018$.
\end{proof}

\subsection{Computational verification}

Among the 922 primes $p \le 10^7$ with $M(p) = -3$:

\begin{center}
\begin{tabular}{lcc}
\toprule
Quantity & Observed & Predicted \\
\midrule
Circular mean of $\gamma_1 \log p$ for $T > 0$ & 5.28 & 5.208 \\
Resultant length $R$ for $T > 0$ & 0.77 & --- \\
Resultant length $R$ for $T < 0$ & 0.13 & --- \\
Fraction $T > 0$ in $[10^6, 10^7)$ & 0.462 & 0.47 \\
\bottomrule
\end{tabular}
\end{center}

The phase match to $0.07$ radians ($1.1\%$) is the sharpest confirmation.
The resultant length $R = 0.77$ for the $T > 0$ primes indicates
extremely strong phase concentration, while $R = 0.13$ for $T < 0$
indicates near-uniform distribution.

%% FIGURE: Phase-lock scatter plot.  Polar plot of gamma_1 * log(p) mod 2*pi
%% for each of the 922 primes.  T > 0 primes shown as filled dots (concentrated
%% in the arc [4.2, 5.8]), T < 0 primes as open circles (uniformly distributed).
%% Mark the predicted phase 5.208 with a radial line.

\subsection{Extreme clustering}

The 247 primes with $T(N) > 0$ form only 16 runs, with lengths up to
81 consecutive $M(p) = -3$ primes.  This reflects the slow
$\gamma_1$-oscillation: when $\gamma_1 \log N$ is in the favorable
phase window $[4.2, 5.8]$, all nearby primes share $T > 0$.  The
phase window width in $N$ is $N \cdot 1.6/\gamma_1 \approx 0.11N$,
explaining why runs of dozens of primes are typical.

\subsection{Comparison with classical Chebyshev bias}

In the classical prime race $\pi(x;4,3) > \pi(x;4,1)$,
Rubinstein--Sarnak showed the logarithmic density of the bias is
$0.9959\ldots$, driven by a constant offset $1/(2\sqrt{x})$ from
the exceptional prime $p = 2$.  Our problem differs:
\begin{itemize}
  \item The offset $-2\log N/\sqrt{N}$ \emph{vanishes} as $N \to \infty$.
  \item The limiting density is $1/2$ (no permanent bias), unlike
    $0.9959$.
  \item The bias is a \emph{Farey regularity} question, not a
    prime-counting question.
\end{itemize}



%% ================================================================
%% SECTION 7: L-FUNCTION SPECTRAL DECOMPOSITION
%% ================================================================
\section{$L$-Function Spectral Decomposition}\label{sec:l-function}

\subsection{Twisted summatory functions}

For a Dirichlet character $\chi \bmod q$, define the twisted Mertens
function $M_\chi(x) = \sum_{k \le x} \mu(k)\chi(k)$ and the twisted
summatory function
\begin{equation}\label{eq:T-chi}
  T_\chi(N) = \sum_{m=2}^{N} \frac{M_\chi(\lfloor N/m \rfloor)}{m}.
\end{equation}
When $\chi = \chi_0$ (the trivial character mod~1), this
reduces to the original~$T(N)$.

\begin{theorem}[Spectral decomposition principle]\label{thm:spectral-decomp}
Per-step Farey discrepancy admits a spectral decomposition indexed by
Dirichlet characters, where each component $T_\chi(N)$ oscillates at
frequencies determined by the nontrivial zeros of~$L(s,\chi)$.

More precisely, the Perron integral representation
\begin{equation}\label{eq:perron-chi}
  T_\chi(N) + M_\chi(N)
  = \frac{1}{2\pi i}\int_{(c)} N^s \cdot
  \frac{L(s{+}1,\chi)}{s\,L(s,\chi)}\,ds
\end{equation}
has poles at the nontrivial zeros of $L(s,\chi)$, with residues
\begin{equation}\label{eq:residue-chi}
  c_{\rho,\chi} = \frac{L(\rho{+}1,\chi)}{\rho\, L'(\rho,\chi)}
  \cdot N^\rho.
\end{equation}
\end{theorem}

\subsection{Computational evidence}

\begin{center}
\begin{tabular}{lccc}
\toprule
Test & $R(T{>}0)$ & $\sigma$ & Verdict \\
\midrule
$T_{\mathrm{plain}}$ at $\gamma_1(\zeta) = 14.13$ & 0.849 & 4.1 & Strong \\
$T_{\chi_{-4}}$ at $\gamma_1(L,\chi_{-4}) = 6.02$ & 0.156 & 3.2 & Clear \\
$T_{\chi_3}$ at $\gamma_1(L,\chi_3) = 8.04$ & 0.276 & 5.1 & Very strong \\
$T_{\mathrm{plain}}$ at $\gamma_2(\zeta) = 21.02$ & 0.547 & 2.6 & Moderate \\
\midrule
\multicolumn{4}{c}{\textbf{Controls}} \\
$T_{\chi_{-4}}$ at $\gamma_1(\zeta) = 14.13$ & 0.033 & 0.7 & Clean null \\
$T_{\chi_3}$ at $\gamma_1(\zeta) = 14.13$ & 0.046 & 0.9 & Clean null \\
\bottomrule
\end{tabular}
\end{center}

The controls confirm selectivity: $T_\chi$ phase-locks at zeros of
$L(s,\chi)$, \emph{not} at zeta zeros.  The three $T$-functions
$T_{\mathrm{plain}}$, $T_{\chi_{-4}}$, $T_{\chi_3}$ are nearly
uncorrelated ($|r| < 0.15$), confirming they carry independent
arithmetic information.

%% FIGURE: Four-panel L-function phase-lock comparison.
%% Top-left: T_plain at gamma_1(zeta) -- strong lock (R=0.85).
%% Top-right: T_chi4 at gamma_1(L,chi4) -- clear lock (R=0.16/0.54).
%% Bottom-left: T_chi3 at gamma_1(L,chi3) -- very strong lock (R=0.28/0.64).
%% Bottom-right: T_chi4 at gamma_1(zeta) -- clean null (R=0.03).

\subsection{Extension to higher moduli}

Testing all primes $\le 50{,}000$ (5{,}133 primes) with characters
mod~5, 7, and~8:

\begin{center}
\begin{tabular}{lcc}
\toprule
Character & $\max(\sigma)$ & Verdict \\
\midrule
$\chi_5$ (quadratic) at own zero & 10.3 & Strong \\
$\chi_7$ (quadratic) at own zero & 12.8 & Strong \\
$\chi_8$ (real primitive) at own zero & 12.4 & Strong \\
\bottomrule
\end{tabular}
\end{center}

\textbf{Every character tested shows phase-lock at its own
$L$-function zero}, with signal strengths $10$--$13\times$ random
expectation.  Higher-order characters (cubic mod~7, quartic mod~5)
also phase-lock, confirming universality across character orders.

\subsection{Non-primitive characters}

For $\chi \bmod 12$ induced from $\chi_3$, the phase-lock occurs at
zeros of $L(s, \chi_3)$ (the primitive character), with
$\sigma = 20.1\times$ for all primes---\emph{stronger} than the
primitive signal.  This is predicted by the factorization
$L(s, \chi) = L(s, \chi^*) \prod_{p \mid q,\, p \nmid q^*}
(1 - \chi^*(p)/p^s)$: the zeros of the induced $L$-function are
exactly those of the primitive $L$-function, plus finitely many from
the Euler product correction.

\subsection{Conjugate symmetry}

$T_{\bar\chi}(N) = \overline{T_\chi(N)}$ exactly, since $\mu$ is
real-valued: $M_{\bar\chi}(x) = \overline{M_\chi(x)}$.  Thus
conjugate characters carry the same information.



%% ================================================================
%% SECTION 8: THE ORDERING PRINCIPLE
%% ================================================================
\section{The Ordering Principle: Farey vs.\ Stern--Brocot}
\label{sec:ordering}

\subsection{The Stern--Brocot tree}

The Stern--Brocot tree generates all positive rationals via mediant
insertion.  Restricted to $[0,1]$, level~$n$ has $2^n$ new fractions,
and the cumulative set $S_n$ has $2^n + 1$ elements.  The per-level
discrepancy is
$\Delta W_{\mathrm{SB}}(n) = W(S_n) - W(S_{n-1})$.

\subsection{Results: no spectral structure}

\begin{center}
\begin{tabular}{lcc}
\toprule
Property & Stern--Brocot & Farey \\
\midrule
$\Delta W$ behavior & Monotone growth & Oscillation \\
Sign changes (first 21 steps) & 1 & 14 \\
Growth & $\Delta W \sim C \cdot 2^n$ & $\Delta W \sim O(1/N)$ \\
Spectral content & None (pure geometric) & Locked to $\zeta$ zeros \\
Governing function & Minkowski $?(x)$ & Riemann $\zeta$ / Mertens \\
Ordering principle & Tree depth (CF structure) & Denominator bound \\
\bottomrule
\end{tabular}
\end{center}

The Stern--Brocot per-level discrepancy satisfies
$\Delta W_{\mathrm{SB}}(n) / 2^{n-1} \to 0.00789\ldots$
(the integral of $(?(x) - x)^2$ against Lebesgue measure,
where~$?$ is Minkowski's question mark function), with ratio
converging to exactly~$2$ and \emph{zero oscillatory content}.

\subsection{The Calkin--Wilf tree is identical}

The cumulative sets $S_n$ are the same for both trees at every level
(verified exactly to level~11, numerically to level~22).  The
discrepancies are identical:
$\Delta W_{\mathrm{CW}} = \Delta W_{\mathrm{SB}}$ exactly.

\subsection{Interpretation}

The absence of spectral structure in Stern--Brocot discrepancy confirms
that the zeta-zero oscillation in Farey $\Delta W$ is a consequence of
the \emph{arithmetic ordering by denominator}, not an intrinsic property
of the rational numbers.  The same rationals, ordered by tree depth,
produce qualitatively different dynamics:
\begin{itemize}
  \item \textbf{Farey (denominator ordering):} Step-by-step inclusion
    introduces number-theoretic correlations (Euler totient, M\"obius
    function) that encode zeta zeros in the oscillation of $\Delta W$.
  \item \textbf{Stern--Brocot (tree ordering):} Mediant-based insertion
    has a purely geometric doubling structure.  The discrepancy is
    governed by Minkowski's~$?(x)$, a singular function with no
    connection to $\zeta$.
\end{itemize}

\emph{Different orderings of the same rationals reveal different
spectral data---or none at all.  The Farey ordering by denominator is
special because it engages multiplicative number theory, making zeta
zeros visible.}

%% FIGURE: Side-by-side comparison.
%% Left panel: Farey DeltaW(N) for N = 2, ..., 200 showing oscillation.
%% Right panel: Stern-Brocot DeltaW(n) for n = 1, ..., 22 showing
%% monotone exponential growth.  Mark the ratio 2.000 on the right panel.



%% ================================================================
%% SECTION 9: UNCONDITIONAL RESULTS
%% ================================================================
\section{Unconditional Results}\label{sec:unconditional}

The results of Sections~\ref{sec:perron}--\ref{sec:phaselock} are
conditional on GRH (and LI for the density statement).  We now give
what can be proved without any unproven hypothesis.

\begin{theorem}[Unconditional oscillation]\label{thm:unconditional}
$T(N) = \Omega_\pm(\sqrt{N})$.  That is:
\[
  \limsup_{N \to \infty} \frac{T(N)}{\sqrt{N}} > 0, \qquad
  \liminf_{N \to \infty} \frac{T(N)}{\sqrt{N}} < 0.
\]
\end{theorem}

\begin{proof}
Adapting the method of Ingham~\cite{Ingham1942}.  The generating
function of $T(N)$ (excluding the $m = 1$ term) corresponds to the
Dirichlet series $[\zeta(s+1)/(s\,\zeta(s))] - [1/\zeta(s)]
= [\zeta(s+1)/s - 1]/\zeta(s)$.

Near a zero $\rho = 1/2 + i\gamma$ of $\zeta$ on the critical line
(guaranteed to exist by Hardy's theorem~\cite{Hardy1914}), the
factor $[\zeta(\rho+1)/\rho - 1]$ is a nonzero constant, so this
Dirichlet series has a singularity at $s = \rho$.

Suppose for contradiction that $T(N) < C\sqrt{N}$ for all~$N$.  Then
the Dirichlet series converges for $\operatorname{Re}(s) > 1/2$,
forcing analyticity there.  But zeta has zeros on
$\operatorname{Re}(s) = 1/2$ (a positive proportion, by
Selberg~\cite{Selberg1942}), giving singularities at these points.
Contradiction.

The $\Omega_-$ direction follows by the same argument applied
to~$-T(N)$.
\end{proof}

\begin{corollary}\label{cor:infinitely-many}
Among primes $p$ with $M(p) = -3$:
\begin{enumerate}[\upshape(a)]
  \item Infinitely many have $\Delta W(p) > 0$
    \textup{(from the $\Omega_+$ result)};
  \item Infinitely many have $\Delta W(p) < 0$
    \textup{(from the Sign Theorem to $10^5$)}.
\end{enumerate}
The lower logarithmic density of each set is positive.
\end{corollary}

\begin{theorem}[Effective upper bound]\label{thm:effective-bound}
There exists an absolute constant $c > 0$ \textup{(ineffective)} such
that
\[
  |T(N)| \le N \cdot \exp\!\bigl(-c\,(\log N)^{3/5}
  (\log\log N)^{-1/5}\bigr)
\]
for all sufficiently large~$N$.
\end{theorem}

\begin{proof}
Apply the Walfisz bound $|M(x)| \le x\,\exp(-c(\log x)^{3/5}
(\log\log x)^{-1/5})$ to each term of
$T(N) = \sum_{m=2}^{N} M(\lfloor N/m \rfloor)/m$, splitting at
$m = \sqrt{N}$.  The low-$m$ terms contribute at most
$O(N\,\exp(-c'(\log N)^{3/5}/(\log\log N)^{1/5}))$; the high-$m$
terms contribute $O(\sqrt{N}\log N)$, which is dominated.
\end{proof}

\begin{remark}
The gap between the oscillation result ($\Omega_\pm(\sqrt{N})$)
and the upper bound ($O(N/\exp(\cdots))$) is enormous.  Under GRH,
$|T(N)| = O(\sqrt{N}\log^2 N)$, nearly matching the oscillation.
Closing this gap unconditionally would constitute major progress on
zeta-zero distribution.
\end{remark}

\subsection{Partial explicit formula without GRH}

\begin{proposition}[Unconditional]\label{prop:partial-explicit}
For any $\sigma \in (1/2, 1)$ and $T > 0$:
\[
  F(N) = -2\log N + C_0
  + \!\!\!\sum_{\substack{\rho :\, \operatorname{Re}(\rho) > \sigma \\
  |\operatorname{Im}(\rho)| < T}} \!\!\! c_\rho\,N^\rho
  + O\!\left(\frac{N^\sigma \log^2 T}{T}\right)
  + O(N\,e^{-c\sqrt{\log N}}).
\]
\end{proposition}

Using the Ingham zero-density estimate, the contribution from
off-line zeros (if any exist) is at most $O(N^{1/2+\varepsilon})$,
so the on-line zeros dominate regardless.



%% ================================================================
%% SECTION 10: MULTI-SCALE STRUCTURE
%% ================================================================
\section{Multi-Scale Decomposition of $T(N)$}\label{sec:multiscale}

\subsection{Scale bands}

Decompose $T(N) = T_{\mathrm{low}} + T_{\mathrm{mid}} + T_{\mathrm{high}}$
where:
\begin{align*}
  T_{\mathrm{low}}(N) &= \sum_{m=2}^{10}
    \frac{M(\lfloor N/m \rfloor)}{m}
    & &\text{(large scales: $N/2$ to $N/10$),} \\
  T_{\mathrm{mid}}(N) &= \sum_{m=11}^{\sqrt{N}}
    \frac{M(\lfloor N/m \rfloor)}{m}
    & &\text{(medium scales),} \\
  T_{\mathrm{high}}(N) &= \sum_{m > \sqrt{N}}^{N}
    \frac{M(\lfloor N/m \rfloor)}{m}
    & &\text{(small scales).}
\end{align*}

Among 922 primes $p \le 10^7$ with $M(p) = -3$:

\begin{center}
\begin{tabular}{lrrrr}
\toprule
Band & Mean & Std & Correlation with $T$ & \% of $\Var(T)$ \\
\midrule
$T_{\mathrm{low}}$ & $+5.4$ & 38.0 & 0.980 & 83.6\% \\
$T_{\mathrm{mid}}$ & $-6.7$ & 8.3 & 0.570 & 4.0\% \\
$T_{\mathrm{high}}$ & $-9.8$ & 2.6 & $-0.147$ & 0.4\% \\
\bottomrule
\end{tabular}
\end{center}

The large-scale band alone accounts for 83.6\% of the total variance.
The three leading terms $M(N/2)/2 + M(N/3)/3 + M(N/5)/5$ capture
$R^2 = 0.959$.

%% FIGURE: Multi-scale decomposition.  Three-panel plot showing
%% T_low, T_mid, T_high as functions of the prime index.
%% T_low dominates and oscillates wildly; T_mid is small and negative;
%% T_high is a near-constant negative offset.

\subsection{Multi-zero phase structure}

Not only $\gamma_1$, but all five leading zeta zeros show highly
significant phase-locking for $T > 0$ primes:

\begin{center}
\begin{tabular}{ccccl}
\toprule
Zero & $\gamma_k$ & $R(T{>}0)$ & Rayleigh $p$-value & Verdict \\
\midrule
1 & 14.135 & 0.804 & $8.7 \times 10^{-70}$ & Dominant \\
2 & 21.022 & 0.605 & $8.4 \times 10^{-40}$ & Strong \\
3 & 25.011 & 0.462 & $1.6 \times 10^{-23}$ & Clear \\
4 & 30.425 & 0.610 & $1.7 \times 10^{-40}$ & Anomalously strong \\
5 & 32.935 & 0.400 & $8.9 \times 10^{-18}$ & Clear \\
\bottomrule
\end{tabular}
\end{center}

The anomalous strength of $\gamma_4$ ($R = 0.61$, nearly as strong as
$\gamma_2$) may reflect second-harmonic coupling, since
$\gamma_4/\gamma_1 = 2.15$.

The joint $(\gamma_1, \gamma_2)$ phase space reveals a resonance
region: the cell $[\gamma_1 \in (270^\circ, 360^\circ)] \times
[\gamma_2 \in (90^\circ, 180^\circ)]$ contains 182 primes, of
which 74\% have $T > 0$.  Outside this region, $T > 0$ is rare or
impossible.



%% ================================================================
%% SECTION 11: OPEN PROBLEMS
%% ================================================================
\section{Open Problems}\label{sec:open}

\begin{enumerate}
  \item \textbf{Unconditional density bounds.}
    Prove that the density of $\Delta W > 0$ among $M(p) = -3$ primes
    lies in some explicit interval $[\varepsilon, 1-\varepsilon]$ with
    effective~$\varepsilon > 0$.  The current unconditional result gives
    positive lower logarithmic density but with an ineffective (and
    presumably tiny) constant.

  \item \textbf{Universal rho stability.}
    Prove analytically that $\rho \to -3.88\ldots$ as $p \to \infty$.
    This would reduce the sign question for $\Delta W$ entirely to the
    behavior of $T(N)$.

  \item \textbf{Cross-control anomaly.}
    Explain why $T_{\chi_7}$ shows a $19.4\times$ signal at the
    \emph{zeta} zero when using all primes (but not when restricted to
    $M(p) = -3$).  Is this genuine cross-talk between $L$-functions,
    or a sampling artifact?

  \item \textbf{Higher-degree $L$-functions.}
    Compute $T_f(N)$ for the Ramanujan $\tau$-function and verify
    phase-lock at zeros of $L(s, \Delta)$.  The Perron mechanism
    is universal, but the ``M\"obius inverse'' coefficients of a
    modular form $L$-function have more complex structure.

  \item \textbf{Root number signature.}
    For characters $\chi$ with root number $\varepsilon(\chi) = -1$
    (forcing $L(1/2, \chi) = 0$), the zero at $s = 1/2$ contributes
    a non-oscillatory constant to $T_\chi$ via the Perron residue.
    Detect this as a persistent bias in $T_\chi(N)$.

  \item \textbf{Anomalous $\gamma_4$ strength.}
    The fourth zeta zero ($\gamma_4 = 30.43$) shows phase-lock
    $R = 0.61$, nearly as strong as $\gamma_2$.  Is this
    second-harmonic coupling ($\gamma_4 \approx 2\gamma_1$), or does
    it reflect deeper structure in the zero distribution?

  \item \textbf{Stern--Brocot spectral structure.}
    Does the Stern--Brocot per-level discrepancy connect to the
    spectrum of the Gauss map transfer operator?  The governing
    function $?(x)$ arises from the continued fraction algorithm,
    whose spectral theory involves the Ruelle transfer operator.

  \item \textbf{Arithmetic geometry interpretation.}
    When fractions $k/p$ enter $\F_p$, they correspond to $p$-torsion
    points on modular curves.  Does $\Delta W(p)$ have a geometric
    interpretation as a regularity measure for torsion distribution
    among cusps?

  \item \textbf{Operator-theoretic formulation.}
    Is there an operator $A_N$ on $L^2(\mathrm{PSL}(2,\Z)
    \backslash \mathbb{H})$ whose trace equals $T(N)$?  If so, the
    Selberg trace formula would decompose $T(N)$ into spectral and
    geometric contributions, embedding our Perron integral into the
    arithmetic quantum chaos framework.

  \item \textbf{Effective computation of the Chebyshev bias.}
    Compute the characteristic function
    $\varphi(t) = \prod_k J_0(2|c_k|t)$ to sufficient precision to
    determine the bias at $N = 10^{10}$, $10^{12}$, and test whether
    the density of $T > 0$ converges to~$1/2$ as predicted.
\end{enumerate}



%% ================================================================
%% SUMMARY OF LOGICAL DEPENDENCIES
%% ================================================================
\section*{Logical dependencies}

\begin{center}
\begin{tabular}{lll}
\toprule
Result & Status & Depends on \\
\midrule
Four-term decomposition (Thm.~\ref{thm:four-term}) & Unconditional & Algebra \\
Regression identity (Thm.~\ref{thm:regression}) & Unconditional & Permutation symmetry \\
Alpha formula (Thm.~\ref{thm:alpha-formula}) & Unconditional & Classical Farey asymptotics \\
Sign Theorem to $10^5$ (Thm.~\ref{thm:sign}) & Computational & Exact arithmetic \\
Counterexample at $243{,}799$ (Prop.~\ref{prop:counterexample}) & Computational & Exact arithmetic \\
Perron integral for $T(N)$ (Thm.~\ref{thm:perron}) & \textbf{GRH} & Perron formula \\
Phase-lock, density $1/2$ (Thm.~\ref{thm:phase-lock}) & \textbf{GRH + LI} & Rubinstein--Sarnak \\
Unconditional oscillation (Thm.~\ref{thm:unconditional}) & Unconditional & Ingham 1942 \\
$L$-function decomposition (Thm.~\ref{thm:spectral-decomp}) & \textbf{GRH} & Perron for $L(s,\chi)$ \\
Stern--Brocot null test & Computational & Direct enumeration \\
\bottomrule
\end{tabular}
\end{center}



%% ================================================================
%% REFERENCES
%% ================================================================
\begin{thebibliography}{99}

\bibitem{Franel1924}
J.~Franel, Les suites de Farey et le probl\`eme des nombres
premiers, \emph{G\"ott.\ Nachr.} (1924), 198--201.

\bibitem{Landau1924}
E.~Landau, Bemerkungen zu der vorstehenden Abhandlung von Herrn
Franel, \emph{G\"ott.\ Nachr.} (1924), 202--206.

\bibitem{HardyWright2008}
G.\,H.~Hardy and E.\,M.~Wright, \emph{An Introduction to the Theory
of Numbers}, 6th ed., Oxford University Press, 2008.

\bibitem{Edwards1974}
H.\,M.~Edwards, \emph{Riemann's Zeta Function}, Academic Press, 1974.

\bibitem{RubinsteinSarnak1994}
M.~Rubinstein and P.~Sarnak, Chebyshev's bias,
\emph{Experiment.\ Math.} \textbf{3} (1994), 173--197.

\bibitem{Ingham1942}
A.\,E.~Ingham, On two conjectures in the theory of numbers,
\emph{Amer.\ J.\ Math.} \textbf{64} (1942), 313--319.

\bibitem{Hardy1914}
G.\,H.~Hardy, Sur les z\'eros de la fonction $\zeta(s)$ de Riemann,
\emph{C.\,R.\ Acad.\ Sci.\ Paris} \textbf{158} (1914), 1012--1014.

\bibitem{Selberg1942}
A.~Selberg, On the zeros of Riemann's zeta-function,
\emph{Skr.\ Norske Vid.\ Akad.\ Oslo} (1942), No.~10.

\bibitem{Titchmarsh1986}
E.\,C.~Titchmarsh, \emph{The Theory of the Riemann Zeta-Function},
2nd ed.\ (revised by D.\,R.~Heath-Brown), Oxford University Press, 1986.

\bibitem{Walfisz1963}
A.~Walfisz, \emph{Weylsche Exponentialsummen in der neueren
Zahlentheorie}, VEB Deutscher Verlag der Wissenschaften, 1963.

\bibitem{KarvonenZhigljavsky2025}
T.~Karvonen and A.~Zhigljavsky, Maximum mean discrepancy of Farey
sequences and the Riemann Hypothesis,
\emph{Acta Math.\ Hungar.} (2025).

\bibitem{Huxley1971}
M.\,N.~Huxley, The distribution of Farey points, I,
\emph{Acta Arith.} \textbf{18} (1971), 281--287.

\bibitem{Devin2017}
L.~Devin, Chebyshev's bias for analytic $L$-functions,
\emph{Math.\ Proc.\ Cambridge Philos.\ Soc.} (2020).
arXiv:1706.06394.

\bibitem{ElMarraki1995}
M.~El Marraki, Fonction sommatoire de la fonction de M\"obius, 2.
Majorations asymptotiques effectives et fortes,
\emph{J.\ Th\'eor.\ Nombres Bordeaux} \textbf{7} (1995), 407--433.

\bibitem{MarklofStrombergsson2003}
J.~Marklof and A.~Str\"ombergsson, The distribution of free path
lengths in the periodic Lorentz gas and related lattice point
problems, \emph{Ann.\ of Math.} \textbf{172} (2010), 1949--2033.

\bibitem{AlkanLedoanZaharescu}
E.~Alkan, A.\,H.~Ledoan, and A.~Zaharescu, On Dirichlet $L$-functions
and the index of visible points,
\emph{Illinois J.\ Math.} \textbf{51} (2007), 455--477.

\end{thebibliography}

\end{document}
